\section{Desarrollo}
Esta seccion debe explicar paso a paso como se construyo y configuro la simulacion.

\subsection{Creación del modelo}
Describe la geometria creada o importada en ANSYS. Incluye dimensiones, unidades y simplificaciones realizadas.

\captura[0.62]{figures/captura_placeholder.png}{Ejemplo de captura del modelo en ANSYS. Reemplaza esta imagen por tu captura real.}

\subsection{Configuración de unidades y dimensiones}

\begin{table}[H]
\centering
\caption{Parámetros principales utilizados en la simulación.}
\label{tab:parametros}
\begin{tabularx}{\linewidth}{>{\bfseries}l X c}
\toprule
Parametro & Descripcion & Valor \\
\midrule
Material & Material asignado al dominio principal & Cobre / FR4 / Aire \\
Frecuencia & Rango de barrido configurado & \SI{1}{GHz}--\SI{10}{GHz} \\
Mallado & Tipo de malla empleada & Adaptativa \\
Solver & Metodo de solucion seleccionado & Modal / Transient / Static \\
\bottomrule
\end{tabularx}
\end{table}

\subsection{Aplicación de simetría}
Explica las fronteras asignadas, puertos, cargas, restricciones, fuentes o condiciones iniciales.

\subsection{Condiciones de frontera}
Indica el tipo de analisis realizado, numero de iteraciones, criterios de convergencia, barridos parametricos y configuracion del solver.
\subsection{Puertos de excitación}
\subsection{Configuración de simulación}
